\hypertarget{the-c-compilation-execution-model}{%
\chapter{THE C++ COMPILATION \& EXECUTION
MODEL}\label{the-c-compilation-execution-model}}

To truly understand C++, you must understand how your code transforms
from text to a running process. This section demystifies the ``black
box'' of the compiler.

\hypertarget{the-build-pipeline-from-source-to-binary}{%
\subsection{1.5.1 The Build Pipeline: From Source to
Binary}\label{the-build-pipeline-from-source-to-binary}}

The ``compilation'' process actually consists of four distinct stages:

\begin{enumerate}
\def\labelenumi{\arabic{enumi}.}
\tightlist
\item
  \textbf{Preprocessing}: Text substitution.

  \begin{itemize}
  \tightlist
  \item
    Removes comments.
  \item
    Expands macros (\passthrough{\lstinline!\#define!}).
  \item
    Includes headers (\passthrough{\lstinline!\#include!}) recursively.
  \item
    Handles conditionals (\passthrough{\lstinline!\#ifdef!}).
  \item
    \emph{Output}: A single ``Translation Unit'' (pure C++ source code).
  \end{itemize}
\item
  \textbf{Compilation}: Syntax to Assembly.

  \begin{itemize}
  \tightlist
  \item
    \textbf{Lexical Analysis}: Tokenizes source (e.g.,
    \passthrough{\lstinline!int!}, \passthrough{\lstinline!main!},
    \passthrough{\lstinline!\{!}).
  \item
    \textbf{Parsing}: Builds Abstract Syntax Tree (AST).
  \item
    \textbf{Semantic Analysis}: Type checking, overload resolution.
  \item
    \textbf{Optimization}: Dead code elimination, loop unrolling,
    inlining (O1, O2, O3).
  \item
    \textbf{Code Generation}: Generates assembly code for the target
    architecture (x86\_64, ARM64).
  \item
    \emph{Output}: Assembly file (\passthrough{\lstinline!.s!} or
    \passthrough{\lstinline!.asm!}).
  \end{itemize}
\item
  \textbf{Assembly}: Assembly to Machine Code.

  \begin{itemize}
  \tightlist
  \item
    Translates mnemonics (\passthrough{\lstinline!MOV!},
    \passthrough{\lstinline!ADD!}) to opcodes
    (\passthrough{\lstinline!0x89!}, \passthrough{\lstinline!0x01!}).
  \item
    \emph{Output}: Object file (\passthrough{\lstinline!.o!} or
    \passthrough{\lstinline!.obj!}). Contains machine code but with
    \emph{unresolved symbols}.
  \end{itemize}
\item
  \textbf{Linking}: Combining Object Files.

  \begin{itemize}
  \tightlist
  \item
    Combines multiple \passthrough{\lstinline!.o!} files and static
    libraries
    (\passthrough{\lstinline!.a!}/\passthrough{\lstinline!.lib!}).
  \item
    Resolves symbols: Matches function calls in one TU to definitions in
    another.
  \item
    Relocates addresses: Adjusts internal pointers.
  \item
    \emph{Output}: Executable (\passthrough{\lstinline!.exe!} or
    ELF/Mach-O) or Shared Library
    (\passthrough{\lstinline!.so!}/\passthrough{\lstinline!.dll!}).
  \end{itemize}
\end{enumerate}

\hypertarget{translation-units-tu-linkage}{%
\subsection{1.5.2 Translation Units (TU) \&
Linkage}\label{translation-units-tu-linkage}}

A \textbf{Translation Unit} is the input to the compiler after
preprocessing. It is the fundamental unit of compilation.

\hypertarget{linkage-types}{%
\subsubsection{Linkage Types}\label{linkage-types}}

Linkage determines if a name (variable/function) is visible to other
TUs.

\begin{enumerate}
\def\labelenumi{\arabic{enumi}.}
\tightlist
\item
  \textbf{No Linkage}:

  \begin{itemize}
  \tightlist
  \item
    Visible only within the current scope (block).
  \item
    Example: Local variables.
  \end{itemize}
\item
  \textbf{Internal Linkage}:

  \begin{itemize}
  \tightlist
  \item
    Visible only within the current Translation Unit.
  \item
    Hidden from the linker.
  \item
    Examples: \passthrough{\lstinline!static!} global variables,
    \passthrough{\lstinline!static!} free functions,
    \passthrough{\lstinline!const!} globals (by default).
  \item
    \emph{Use Case}: Private helper functions.
  \end{itemize}
\item
  \textbf{External Linkage}:

  \begin{itemize}
  \tightlist
  \item
    Visible to the linker and other TUs.
  \item
    Examples: Global variables, non-static functions,
    \passthrough{\lstinline!extern!} variables.
  \item
    \emph{Danger}: Requires strict ODR compliance.
  \end{itemize}
\end{enumerate}

\begin{lstlisting}[language={C++}]
// TU1.cpp
static int internal_var = 10; // Internal Linkage
int global_var = 20;          // External Linkage

// TU2.cpp
extern int global_var;        // Declares existence of external symbol
// extern int internal_var;   // ERROR: Linker error (symbol not found)
\end{lstlisting}

\hypertarget{the-one-definition-rule-odr}{%
\subsection{1.5.3 The One Definition Rule
(ODR)}\label{the-one-definition-rule-odr}}

The ODR is the most important rule in C++ linking.

\begin{enumerate}
\def\labelenumi{\arabic{enumi}.}
\tightlist
\item
  \textbf{ODR for Translation Units}: A function/variable can have only
  \textbf{one definition} in a single TU. (Declarations can be
  repeated).
\item
  \textbf{ODR for Programs}: A non-inline function/variable can have
  only \textbf{one definition} in the entire program.
\item
  \textbf{ODR for Classes/Templates}: Can be defined in multiple TUs
  (via headers), but definitions must be \textbf{token-for-token
  identical}.
\end{enumerate}

\textbf{Violation Example:}

\begin{lstlisting}[language={C++}]
// header.h
int x = 10; // DEFINITION (allocates memory)

// a.cpp
#include "header.h"

// b.cpp
#include "header.h"

// Linker Error: multiple definition of `x`
// Fix: Use 'extern int x;' in header, definition in one .cpp
// Fix (C++17): Use 'inline int x = 10;'
\end{lstlisting}

\hypertarget{process-memory-layout}{%
\subsection{1.5.4 Process Memory Layout}\label{process-memory-layout}}

When your C++ program runs, the OS allocates virtual memory segments:

\begin{enumerate}
\def\labelenumi{\arabic{enumi}.}
\tightlist
\item
  \textbf{Text Segment (.text)}:

  \begin{itemize}
  \tightlist
  \item
    Contains executable machine instructions.
  \item
    Read-Only (prevents self-modifying code exploits).
  \item
    Shared between multiple instances of the app.
  \end{itemize}
\item
  \textbf{Data Segment (.data)}:

  \begin{itemize}
  \tightlist
  \item
    Initialized global and static variables.
  \item
    \passthrough{\lstinline!int g = 10;!} lives here.
  \end{itemize}
\item
  \textbf{BSS Segment (.bss)}:

  \begin{itemize}
  \tightlist
  \item
    Uninitialized global/static variables.
  \item
    \passthrough{\lstinline!int g;!} lives here.
  \item
    Automatically zero-initialized by OS before
    \passthrough{\lstinline!main()!}.
  \end{itemize}
\item
  \textbf{Heap}:

  \begin{itemize}
  \tightlist
  \item
    Dynamic memory (\passthrough{\lstinline!new!},
    \passthrough{\lstinline!malloc!}).
  \item
    Grows upwards (typically).
  \item
    Managed manually or by smart pointers.
  \end{itemize}
\item
  \textbf{Stack}:

  \begin{itemize}
  \tightlist
  \item
    Local variables, function parameters, return addresses.
  \item
    Grows downwards (typically).
  \item
    LIFO (Last-In, First-Out).
  \item
    \emph{Stack Overflow}: Exceeding stack size (e.g., deep recursion).
  \end{itemize}
\end{enumerate}

\hypertarget{program-startup-before-main}{%
\subsection{1.5.5 Program Startup: Before
main()}\label{program-startup-before-main}}

\passthrough{\lstinline!main()!} is NOT the first thing to run.

\begin{enumerate}
\def\labelenumi{\arabic{enumi}.}
\tightlist
\item
  \textbf{OS Loader}: Loads EXE into memory.
\item
  \textbf{Entry Point (\passthrough{\lstinline!\_start!})}: Defined by C
  Runtime (CRT).
\item
  \textbf{Global Initialization}:

  \begin{itemize}
  \tightlist
  \item
    Constructors of global/static objects run \textbf{before}
    \passthrough{\lstinline!main!}.
  \item
    \emph{Static Initialization Order Fiasco}: Order between TUs is
    undefined.
  \end{itemize}
\item
  \textbf{\passthrough{\lstinline!main()!}}: Your code runs.
\item
  \textbf{Global Destruction}:

  \begin{itemize}
  \tightlist
  \item
    Destructors of global/static objects run \textbf{after}
    \passthrough{\lstinline!main!} returns.
  \item
    \passthrough{\lstinline!atexit!} handlers run.
  \end{itemize}
\end{enumerate}

\hypertarget{deep-dive-into-data-representation}{%
\subsection{1.5.6 Deep Dive into Data
Representation}\label{deep-dive-into-data-representation}}

To understand bugs like integer overflow and floating-point inaccuracy,
you must know how data is stored bits-by-bits.

\hypertarget{integer-representation-twos-complement}{%
\subsubsection{Integer Representation (Two's
Complement)}\label{integer-representation-twos-complement}}

Most modern systems use \textbf{Two's Complement} for signed integers. *
\textbf{Positive Numbers}: Standard binary.
\passthrough{\lstinline!0000 0101!} (5). * \textbf{Negative Numbers}:
Invert all bits and add 1. * \passthrough{\lstinline!5!} =
\passthrough{\lstinline!0000 0101!} * Invert:
\passthrough{\lstinline!1111 1010!} * Add 1:
\passthrough{\lstinline!1111 1011!} (-5) * \textbf{Advantage}:
Addition/Subtraction works identically for signed/unsigned. *
\textbf{Range}: \passthrough{\lstinline!-2\^(N-1)!} to
\passthrough{\lstinline!2\^(N-1) - 1!}. (One extra negative number).

\hypertarget{floating-point-ieee-754}{%
\subsubsection{Floating Point (IEEE
754)}\label{floating-point-ieee-754}}

\passthrough{\lstinline!float!} (32-bit) and
\passthrough{\lstinline!double!} (64-bit) follow IEEE 754. *
\textbf{Components}: 1. \textbf{Sign Bit}: 0 (+) or 1 (-). 2.
\textbf{Exponent}: Biased integer (allows very large/small numbers). 3.
\textbf{Mantissa (Significand)}: The precision bits (normalized to
1.xxxxx). * \textbf{Implication}: *
\passthrough{\lstinline"0.1 + 0.2 != 0.3"} because 0.1 cannot be
represented exactly in binary (infinite repeating fraction). *
\textbf{NaN (Not a Number)}: Result of 0/0 or sqrt(-1).
\passthrough{\lstinline"NaN != NaN"} is always true. *
\textbf{Infinity}: Result of 1.0/0.0.

\begin{lstlisting}[language={C++}]
#include <iostream>
#include <cmath>
#include <limits>

int main() {
    float a = 0.1f;
    float b = 0.2f;
    if (a + b == 0.3f) {
        std::cout << "Math works!\n";
    } else {
        std::cout << "Math is broken: " << (a+b) << "\n"; // Prints 0.30000001
    }
    
    // Correct comparison
    if (std::abs((a + b) - 0.3f) < 1e-5) {
        std::cout << "Close enough!\n";
    }
}
\end{lstlisting}

\begin{center}\rule{0.5\linewidth}{0.5pt}\end{center}
